% SPDX-License-Identifier: Apache-2.0
% covers: Pair
\datasheet{Pair}{Two Vreteno cores in step, with everything they say compared and one line that says they stopped agreeing.}

\dsfacts{\code{cpu/vreteno/src/pair.rs}}{\code{vreteno32::pair::Pair}}%
{\code{//docs:vreteno}}

\subsection*{Function}

\code{Pair<IW>} is two \code{Vreteno} cores running one program, with
one line that says they disagreed.

Every input the pair is given goes to both: the wires directly, since a
wire has no handshake, and each of the three channels through a
\code{Tee}, which hands a word to both cores in one cycle or to
neither. Every output is compared: the three channels through a
\code{Check} each, and the three wires through \code{Watch}.

What leaves the pair is the first core's, so a design puts a pair where
it would have put a core and wires nothing differently. What the pair
adds is one input and one output: \code{fault}, the self-test, and
\code{differs}.

\subsection*{Parameters}

\begin{center}\footnotesize
\begin{tabular}{@{}lp{0.7\textwidth}@{}}
\toprule
Parameter & Meaning \\
\midrule
\code{IW} & The width of the AXI identifier the cores use, in bits,
passed to both cores and to the channels between them. The tables use
2, which is the only width in the tree. \\
\bottomrule
\end{tabular}
\end{center}

\dstables{Pair}

\subsection*{Behaviour}

\begin{itemize}
\item The pair holds no register of its own. Every register in it
belongs to a core, a tee, a check or the watch.
\item The three tees put one cycle between a word reaching the pair and
a core being offered it. The delay is the same for both cores, which is
the point of it, and neither core can be offered a word the other has
not had.
\item A disagreement raises \code{differs} in the cycle it happens, on
a channel or on a wire alike, and the line stays up until \code{rst}.
\item Nothing else happens. Neither core is stopped, no channel is
held, no trap is raised, and the first core's word still goes out. The
pair says that there was a disagreement; what to do about it, between a
reset, a trap and a stop, belongs to the system around it.
\item A core that stops offering on a channel stalls that channel
rather than raising the line, since a check compares only words that
move. Such a core shows up on the wires, which are compared every
cycle.
\end{itemize}

\subsection*{Verification}

\begin{itemize}
\item \code{//cpu/vreteno:pair\_test} runs the pair on the lockstep test's
machine, with a tracker, a router, the data memory, the interrupt
controller and the serial port.
\code{two\_cores\_given\_the\_same\_run\_agree} runs eight random
programs and finds the line never up.
\code{a\_core\_held\_in\_reset\_a\_cycle\_longer\_is\_caught} raises
\code{fault} for one cycle and finds it up.
\code{the\_line\_stays\_up\_after\_a\_disagreement} reads it again at the
end of six hundred cycles and finds it still up.
\item The pair is never lowered, so it has no netlist, no
co-simulation and no synthesis. Its two parts from \code{//lib/parts}
are simulated on their own through \code{ex\_redundant}.
\end{itemize}

\subsection*{Used by}

Nothing in the tree but its own test. It is the shape a design takes
when it must not be quietly wrong, and no design here has asked for it
yet.

\subsection*{Limits}

One bit of report, with no record of which comparison fired or when.
The only fault the design can inject is holding the second core in
reset. The two cores are two instances of one Rust type, so a fault in
the design is in both and neither the comparison nor the self-test can
see it; what a pair catches is a fault that strikes one copy. The pair
has no address map and no registers of its own, so a program cannot
read the line; a design brings it out as a wire.
